\section{Background: The Classical Computing Stack and the Model-Native System Stack}
\label{sec:background}

This section first reviews the layered logic of classical computer architecture and the quantitative design principles that underpin it, then surveys the nascent stratification currently taking shape in large language model (LLM) systems. Together, these two perspectives lay the groundwork for the analogy framework developed in subsequent sections.

\subsection{Layered Abstractions in Classical Computer Architecture}

The layered organization of a classical computing system is not an ad hoc design choice but a general engineering strategy for managing complexity~\cite{hennessy2017quantitative}.
The central idea is that each layer depends only on the \emph{interface} exposed by the layer below, never on its \emph{implementation}.
This separation allows any layer to be independently optimized (or even replaced) without disrupting the rest of the stack.
Figure~\ref{fig:classical_stack} provides an overview of the six-layer classical stack and its correspondence to ICA layers.
% ---------------------------------------------------------------------------
% Classical Computing Stack — six-layer abstraction overview
% fig:classical_stack
% ---------------------------------------------------------------------------
\begin{figure}[htbp]
\centering
\providecolor{csL1}{HTML}{2E5E9E}
\providecolor{csL2}{HTML}{3A72B5}
\providecolor{csL3}{HTML}{4A8AC8}
\providecolor{csL4}{HTML}{5EA0D8}
\providecolor{csL5}{HTML}{73B4E4}
\providecolor{csL6}{HTML}{90C8F0}
%
\resizebox{0.85\textwidth}{!}{%
\begin{tikzpicture}[
    >=Latex, font=\small,
    layrow/.style={
        draw=#1!70!black, fill=#1!22, rounded corners=4pt,
        minimum width=14.0cm, minimum height=1.0cm,
        align=center, font=\small\bfseries, text=#1!90!black,
        line width=0.65pt
    },
    annot/.style={font=\scriptsize, text=black!60, align=left, text width=5.8cm},
    classicannot/.style={font=\scriptsize, text=black!55, align=right, text width=5.2cm},
]

\def\gap{1.25}

% ---- Layer rows (bottom to top) ----
\node[layrow=csL1] (hw) at (0, 0*\gap)
    {L1 \enspace ISA / Physical Hardware\enspace
    {\normalfont\scriptsize(x86, ARM, RISC-V; transistors, DRAM, SSD)}};

\node[layrow=csL2] (micro) at (0, 1*\gap)
    {L2 \enspace Microarchitecture\enspace
    {\normalfont\scriptsize(pipeline, superscalar, branch prediction, out-of-order)}};

\node[layrow=csL3] (cache) at (0, 2*\gap)
    {L3 \enspace Cache Hierarchy\enspace
    {\normalfont\scriptsize(L1/L2/L3 SRAM caches; TLB; hardware prefetcher)}};

\node[layrow=csL4] (vm) at (0, 3*\gap)
    {L4 \enspace Virtual Memory\enspace
    {\normalfont\scriptsize(paging, page tables, MMU, demand paging, swap)}};

\node[layrow=csL5] (os) at (0, 4*\gap)
    {L5 \enspace Operating System\enspace
    {\normalfont\scriptsize(process scheduler, IPC, file system, access control)}};

\node[layrow=csL6] (dist) at (0, 5*\gap)
    {L6 \enspace Distributed Systems / Applications\enspace
    {\normalfont\scriptsize(consensus, replication, CAP, Raft, Spanner)}};

% ---- Right annotations: key principle ----
\node[font=\scriptsize, text=black!60, anchor=west] at (7.3, 0*\gap)  {\textit{Stable contract: ISA decouples SW from HW}};
\node[font=\scriptsize, text=black!60, anchor=west] at (7.3, 1*\gap)  {\textit{Amdahl's Law; ILP; power wall $\to$ multicore}};
\node[annot, anchor=west] at (7.3, 2*\gap)  {\textit{Locality principle; hit rate $\to$ speedup}};
\node[annot, anchor=west] at (7.3, 3*\gap)  {\textit{Virtualization; isolation; demand paging}};
\node[annot, anchor=west] at (7.3, 4*\gap)  {\textit{Least privilege; scheduling; abstraction}};
\node[annot, anchor=west] at (7.3, 5*\gap)  {\textit{CAP theorem; consensus; fault tolerance}};

% ---- Left annotations: analogy target in ICA ----
\node[classicannot, anchor=east] at (-7.3, 0*\gap)  {ICA L1\\Physical Execution};
\node[classicannot, anchor=east] at (-7.3, 1*\gap)  {ICA L2\\Inference Serving};
\node[classicannot, anchor=east] at (-7.3, 2*\gap)  {ICA L2--L3\\KV Cache / Context};
\node[classicannot, anchor=east] at (-7.3, 3*\gap)  {ICA L3\\Context Mgmt};
\node[classicannot, anchor=east] at (-7.3, 4*\gap)  {ICA L4--L5\\Interface / Orchestration};
\node[classicannot, anchor=east] at (-7.3, 5*\gap)  {ICA L5--L6\\Multi-Agent / App};


\end{tikzpicture}}%
\caption{The classical computing stack as a six-layer abstraction hierarchy, with corresponding ICA layer mappings on the left and the governing design principle on the right. Each layer depends only on the interface of the layer below, enabling independent evolution. The ICA architecture (Section~\ref{sec:icam}) transplants this layered discipline into the model-native computing domain.}
\label{fig:classical_stack}
\end{figure}

We now review the core concepts of each layer in turn.

\subsubsection{Instruction Set Architecture (ISA): The Contract Between Hardware and Software}

The \textit{Instruction Set Architecture} (ISA) specifies the machine behavior visible to software: the available instructions (arithmetic, memory access, branching, jumping), the register file, how exceptions are raised and handled, and how the address space is organized.
The ISA serves as a stable contract: as long as the processor honors this contract, software above it runs correctly, regardless of the microarchitectural details hidden below.
Intel x86, ARM, and RISC-V embody three distinct design philosophies~\cite{intel_sdm_2026,riscv_spec_2026}: x86 prioritizes backward compatibility (code written four decades ago still executes), ARM optimizes for energy efficiency (a decisive advantage in mobile devices), and RISC-V embraces openness and modularity (users may define custom instruction extensions).
Despite their differences, all three serve the same purpose: providing a stable programming target for the software stack.

\subsubsection{Microarchitecture: Multiple Implementations Under One Contract}

The \textit{microarchitecture} determines how a given ISA is realized in silicon and, consequently, what performance a processor actually delivers.
It answers questions such as: how many clock cycles does an instruction require? How many instructions can execute simultaneously? What is the penalty for a mispredicted branch?
Concretely, a \emph{pipeline} decomposes instruction execution into stages (fetch, decode, execute, memory access, write-back) so that different instructions can overlap across stages.
A \emph{superscalar} design dispatches multiple instructions per cycle to independent execution units.
A \emph{branch predictor} uses historical patterns to speculate on the outcome of conditional branches, thereby avoiding pipeline bubbles.
Both Intel's Core microarchitecture and AMD's Zen microarchitecture implement the same x86 ISA yet deliver different performance characteristics, and this difference is entirely transparent to software.

\subsubsection{Cache Hierarchy: Exploiting Locality to Bridge the Processor--Memory Gap}

The \textit{cache hierarchy} addresses a fundamental speed mismatch between processors and main memory.
A modern CPU core can complete an operation in roughly one nanosecond, whereas a DRAM access requires approximately 100 nanoseconds, forcing the processor to stall for two orders of magnitude longer if every data reference must reach main memory.
Caches exploit the \emph{principle of locality}: \emph{temporal locality} holds that data accessed recently is likely to be accessed again soon, while \emph{spatial locality} holds that data at nearby addresses will likely be accessed in the near future~\cite{drepper2007memory}.
Accordingly, caches store recently and nearby accessed data in faster but smaller SRAM: the L1 cache (typically 32--64\,KB, $\sim$4-cycle latency) sits closest to the core, the L2 cache (256\,KB--1\,MB, $\sim$10 cycles) acts as a second-level buffer, and the L3 cache (several to tens of megabytes, $\sim$40 cycles) is shared among multiple cores.
On a cache miss, data must be fetched from DRAM at a cost of roughly 200 cycles.
This tiered design trades capacity for speed at each level, using SRAM's low latency for hot data and DRAM's density for the working set at large.

\subsubsection{Virtual Memory: The Illusion of Infinite Address Space}

\textit{Virtual memory}, managed by the Memory Management Unit (MMU) and page tables, provides each process with an independent address space that can far exceed physical memory.
The core mechanism is \emph{paging}: the virtual address space is divided into fixed-size \emph{pages} (typically 4\,KB), physical memory into equally sized \emph{page frames}, and a page table records the mapping from virtual page numbers to physical frame numbers.
When a process accesses a virtual page not yet resident in physical memory, a \emph{page fault} is raised; the operating system loads the page from disk into a free frame, updates the mapping, and resumes the process~\cite{ostep_2023}.
Every process thus believes it possesses a complete, private address space, while the operating system transparently manages allocation and reclamation of the underlying physical resource.
Virtual memory also enforces isolation: each process maintains its own page table and cannot address another process's memory.

\subsubsection{Operating System: Scheduling, Protection, and Resource Management}

The \textit{operating system} (OS) is the first software layer above hardware.
It manages compute resources and provides uniform abstractions to applications.
Its core responsibilities include: process scheduling (deciding which process runs next, e.g., Linux's Completely Fair Scheduler~\cite{linux_cfs_2026}), memory management (allocating and reclaiming physical page frames), I/O management (providing uniform file and device interfaces), access control (ensuring processes access only authorized resources), and concurrency support (semaphores, locks, condition variables)~\cite{ostep_2023,xv6_2022}.
The OS design principle is simple: expose clean abstractions upward (processes, files, sockets) while managing complex hardware downward (CPU cores, memory modules, disks, network interfaces).
Applications invoke system calls such as \texttt{open()}, \texttt{read()}, and \texttt{fork()} without needing to know whether the storage device is an SSD or an HDD, or whether the network link runs at 1\,GbE or 10\,GbE.

\subsubsection{Distributed Systems: Scaling Beyond a Single Node}

\textit{Distributed systems} extend the computing stack across multiple nodes organized into a logically unified whole.
The fundamental challenge is maintaining correctness and availability when individual nodes may fail, network latencies are variable, and data must be replicated~\cite{cap2002}.
The Raft consensus algorithm~\cite{raft2014} enables a cluster of nodes to agree on the next command to append to a shared log, tolerating crashes of a minority of participants.
Google Spanner~\cite{spanner2012} provides a globally consistent, distributed database service across multiple data centers worldwide.
The CAP theorem~\cite{cap2002} formalizes an inherent trade-off: during a network partition, a system cannot simultaneously guarantee both consistency (C) and availability (A); a design choice must be made.

\subsubsection{Quantitative Design Principles}

What elevates computer architecture from craft to science is a set of \textbf{quantitative principles} that make performance predictable and compositional.

\paragraph{Amdahl's Law.}
Amdahl's Law establishes a theoretical upper bound on the speedup achievable through parallelization:
\begin{equation}
\label{eq:amdahl}
S = \frac{1}{(1-f) + f/p}
\end{equation}
where $f$ is the fraction of the workload that can be parallelized, $p$ is the number of processors, and $S$ is the resulting speedup.
For instance, if 80\% of a task is parallelizable ($f=0.8$) and four processors are available ($p=4$), the maximum speedup is $S = 1/(0.2 + 0.8/4) = 2.5\times$, far below the ideal $4\times$.
Even as $p \to \infty$, speedup asymptotically approaches $1/(1-f) = 5\times$.
The key insight is that \emph{the serial fraction caps the speedup ceiling}.
Malekar and Zand (2024) adapted Amdahl's Law to the LLM setting, proposing an analytical framework for LLM throughput~\cite{malekar2024amdahlllm}.

\paragraph{The Roofline Model.}
The Roofline model visualizes the performance constraints of an operator on a given hardware platform~\cite{hennessy2017quantitative}.
The attainable performance is bounded by two lines: a horizontal line representing peak compute throughput (FLOP/s) and a sloped line representing memory bandwidth multiplied by arithmetic intensity ($\mathrm{Bandwidth} \times \mathrm{Arithmetic\ Intensity}$).
Their intersection is called the \emph{ridge point}.
Operators with arithmetic intensity below the ridge point are \emph{memory-bandwidth-bound}; those above it are \emph{compute-bound}.
Yuan et al.\ systematically applied the Roofline model to LLM inference, revealing how different optimization strategies trade off bandwidth and compute utilization~\cite{yuan2024llminference}.

\paragraph{Locality Principle.}
Temporal and spatial locality directly inform cache design: temporal locality (recently accessed data will be re-accessed) motivates cache residency, while spatial locality (nearby addresses will be accessed soon) motivates cache lines (typically 64 bytes), hardware prefetching, and replacement policies such as LRU~\cite{drepper2007memory}.

Taken together, these quantitative principles have enabled the computing stack to scale from thousands of floating-point operations per second to over $10^{18}$ operations per second across eight decades, all without requiring programmers to understand the physical details of every layer.

\subsection{Emerging Layers of the LLM System Stack}

LLM systems are currently undergoing a stratification process that closely mirrors the classical computing stack.
We survey each layer from bottom to top, highlighting key developments and the system-level abstractions they implicitly introduce.

\subsubsection{Model Layer: Evolution of the Inference Core}

On the modeling front, the Transformer architecture~\cite{vaswani2017attention} remains the dominant backbone.
Its self-attention mechanism captures long-range dependencies by computing pairwise interactions across all positions in a sequence.
Open-weight models continue to evolve along several axes: RoPE (Rotary Position Embedding)~\cite{su2021roformer} improves positional representations, enabling models to handle longer sequences more effectively; Mixture-of-Experts (MoE)~\cite{switch2022,mixtral2024} scales model capacity through sparse activation (activating only a subset of parameters per inference), avoiding a proportional increase in computation; long-context extensions~\cite{yarn2024,longrope2024} push the effective context window from thousands of tokens to hundreds of thousands or even millions.
In parallel, selective state-space models such as Mamba~\cite{gu2024mamba} open a linear-complexity sequence modeling paradigm that no longer relies on self-attention.

\subsubsection{Inference Layer: Systematic Serving Optimizations}

On the inference front, a growing body of work has transformed LLM serving from a single forward pass into a systems engineering problem.
FlashAttention~\cite{dao2022flashattention,dao2023flashattention2} reformulates attention computation as an I/O-aware problem: by tiling the computation to reduce accesses to high-bandwidth memory (HBM), it lowers attention memory complexity from $O(N^2)$ to $O(N)$.
vLLM's PagedAttention~\cite{kwon2023pagedattention} models KV cache management as a paging problem, dividing the KV cache into fixed-size ``pages'' that are allocated and reclaimed on demand, thereby eliminating the waste of pre-allocating maximum memory per request.
ORCA~\cite{orca2022} introduces continuous batching at the iteration level, allowing new requests to join a running batch during the decode phase of existing requests.
DistServe~\cite{zhong2024distserve} decouples the prefill phase (processing the input prompt) from the decode phase (generating tokens one at a time) onto separate GPUs, avoiding resource contention between the two.
Sarathi-Serve~\cite{agrawal2024sarathi} further refines the throughput--latency trade-off.
Disaggregated serving architectures push this prefill/decode decoupling to cluster scale: Splitwise~\cite{patel2024splitwise} and Mooncake~\cite{qin2024mooncake} separate the two phases across dedicated GPU pools and treat the KV cache as a global, transferable resource, realizing the ``KV cache as a hardware-managed cache'' analogy at hyperscale.
A unifying theme across these systems is the adaptation of classical operating systems techniques (paging, batch scheduling, resource isolation) to LLM inference engines.

\subsubsection{Memory Layer: Reconciling Finite Windows with Persistent State}

At the memory layer, the core challenge is that the context window is finite (e.g., 128K tokens), whereas agents often need access to knowledge and interaction histories that far exceed this limit.
Multiple technical directions are being explored in parallel.
Long-context extensions~\cite{gemini2024,yarn2024} attempt to enlarge the window directly, but face quadratic growth in computation cost with sequence length.
Retrieval-Augmented Generation (RAG)~\cite{lewis2020rag} retrieves relevant passages from an external knowledge base at inference time and injects them into the context, analogous to demand paging in an operating system.
MemGPT~\cite{memgpt2023} explicitly introduces \textit{virtual context management}, partitioning context into a working context (analogous to main memory) and an external storage (analogous to disk), and automatically moving information between the two.
LongMem~\cite{longmem2023} and Generative Agents~\cite{park2023generativeagents} explore alternative forms of long-term memory.
LongMemEval~\cite{longmemeval2024} establishes a benchmark for evaluating memory in long-term interactions.
Letta~\cite{letta_stateful_2026} places stateful agents at the center of its design, enabling agents to retain memory across sessions.
MemOS~\cite{memos2025} goes further by treating memory as a first-class operating-system-level resource, unifying management of textual memory, activation memory, and other representational forms.
MemoryOS~\cite{memoryos2025}, presented at EMNLP 2025, proposes a memory operating system tailored to AI agents.
\emph{Context engineering}~\cite{contextengineer_survey2025} has emerged as a nascent discipline that systematizes this space.
It is defined as the engineering practice of systematically designing, selecting, and optimizing all information presented to an LLM, encompassing retrieval augmentation, memory management, tool-use scaffolding, and multi-turn dialogue optimization.

\subsubsection{Agent Layer: From Single-Turn Inference to Complex Runtimes}

At the agent layer, LLMs have evolved beyond single-turn inference engines into controllers within complex runtime environments.
ReAct~\cite{yao2023react} interleaves reasoning with acting, enabling a model to solve multi-step tasks in a ``think--act--observe'' loop.
Toolformer~\cite{schick2023toolformer} teaches models to invoke external tools (calculators, search engines) autonomously.
SWE-agent~\cite{sweagent2024} and OpenHands~\cite{openhands_paper_2025} are autonomous agents for software engineering that can navigate code repositories, write patches, and run test suites.
Codex~\cite{openai_codex_overview_2026} and Claude Code~\cite{anthropic_claudecode_overview_2026} are production-grade coding agents equipped with OS-kernel-like capabilities: sub-agent scheduling~\cite{openai_codex_overview_2026,anthropic_claudecode_subagents_2026}, sandbox isolation~\cite{openai_codex_sandbox_2026,anthropic_claudecode_sandbox_2026}, and permission control~\cite{openai_codex_approvals_2026,anthropic_claudecode_security_2026}.
Multi-agent collaboration frameworks such as AutoGen~\cite{wu2023autogen} and MetaGPT~\cite{metagpt2023} further organize ensembles of specialized agents into cooperative networks.

\subsection{The Gap: No Unified System Language}

Despite rapid progress across each layer, the LLM systems landscape lacks a unified system language, that is, a shared vocabulary and abstraction framework that connects optimizations across layers.
Without such a framework, research at each layer tends to fragment into isolated collections of engineering tricks.
Consider, for example, \emph{prefix caching} (reusing cached key--value pairs across shared prompt prefixes), \emph{chunked prefill} (breaking long prefill computations into schedulable chunks), \emph{sub-agents} (concurrently executing child agents), \emph{MCP servers} (standardized tool interfaces via the Model Context Protocol), and \emph{persistent memory} (durable storage for agent state).
These concepts appear to belong to entirely different research directions, yet from an architectural standpoint they correspond to cache reuse, scheduling granularity, concurrent execution, I/O interfaces, and tiered storage, all well-studied abstractions in classical systems.

This observation motivates the architectural analogy that runs through this paper.
The analogy does not claim that LLM systems and classical systems are ``fundamentally the same.''
Rather, it places independently developed engineering techniques back into a larger, unified design space, making it possible to compose, predict, and systematically design cross-layer optimizations.

Figure~\ref{fig_dual_stack} illustrates the structural correspondence between the two stacks.
% Preamble:
% \usepackage{tikz}
% \usetikzlibrary{arrows.meta,positioning,decorations.pathreplacing,calc}

\definecolor{classical}{HTML}{3D6FB6}
\definecolor{modelnative}{HTML}{D9792B}
\definecolor{ink}{HTML}{22303C}
\definecolor{softgray}{HTML}{6E7B87}
\definecolor{bgline}{HTML}{D8DEE6}
\definecolor{pillbg}{HTML}{F5F7FA}

% Classical stack colors (bottom -> top)
\definecolor{c1}{HTML}{234C84}
\definecolor{c2}{HTML}{2F5F9E}
\definecolor{c3}{HTML}{3E73B2}
\definecolor{c4}{HTML}{5489C4}
\definecolor{c5}{HTML}{6AA2D8}
\definecolor{c6}{HTML}{86BDE8}

% Model-native stack colors (bottom -> top)
\definecolor{m1}{HTML}{A14A12}
\definecolor{m2}{HTML}{BC5C18}
\definecolor{m3}{HTML}{D36E20}
\definecolor{m4}{HTML}{E38431}
\definecolor{m5}{HTML}{EE9D50}
\definecolor{m6}{HTML}{F4B777}

\newcommand{\subline}[1]{{\color{white!92}\tiny #1}}

\begin{figure}[htbp]
\centering
\resizebox{0.98\linewidth}{!}{%
\begin{tikzpicture}[
    x=1cm,y=1cm,
    >={Stealth[length=2.2mm,width=1.4mm]},
    font=\sffamily,
    box/.style={
        rounded corners=4pt,
        minimum width=5.55cm,
        minimum height=1.02cm,
        align=center,
        text=white,
        inner xsep=5pt,
        inner ysep=4pt,
        line width=0.5pt,
        draw=white!35,
        font=\small
    },
    stacktitle/.style={
        rounded corners=8pt,
        fill=pillbg,
        draw=bgline,
        line width=0.5pt,
        minimum height=0.55cm,
        inner xsep=10pt,
        font=\normalsize\bfseries
    },
    layernum/.style={font=\scriptsize\bfseries,text=softgray},
    rel/.style={-{Stealth[length=2.0mm,width=1.3mm]},draw=bgline,line width=0.9pt},
    maplabel/.style={
        font=\scriptsize\bfseries,
        text=ink,
        fill=white,
        draw=bgline,
        rounded corners=2.2pt,
        inner xsep=4pt,
        inner ysep=1.7pt
    },
    sidetitle/.style={font=\scriptsize\bfseries,align=center}
]

% Section headers
\node[stacktitle,text=classical] at (-3.8,7.25) {Classical Computing Stack};
\node[stacktitle,text=modelnative] at ( 3.8,7.25) {Model-Native Computing Stack (ICAM)};

% Central divider
\draw[bgline, line width=0.8pt] (0,-0.85) -- (0,6.95);

% Column boxes, left
\node[box, fill=c6] (cl6) at (-3.8, 6.05) {\textbf{Distributed Systems}\\[-1pt]\subline{MPI $\cdot$ RPC $\cdot$ Load Balancing $\cdot$ Consensus}};
\node[box, fill=c5] (cl5) at (-3.8, 4.85) {\textbf{Operating System}\\[-1pt]\subline{Scheduling $\cdot$ Syscalls $\cdot$ IPC $\cdot$ Virtualization $\cdot$ File System}};
\node[box, fill=c4] (cl4) at (-3.8, 3.65) {\textbf{Virtual Memory}\\[-1pt]\subline{Page Tables $\cdot$ TLB $\cdot$ Demand Paging $\cdot$ Working Set $\cdot$ Thrashing}};
\node[box, fill=c3] (cl3) at (-3.8, 2.45) {\textbf{Cache Hierarchy}\\[-1pt]\subline{L1/L2/L3 $\cdot$ Write-Back $\cdot$ Cache Coherence}};
\node[box, fill=c2] (cl2) at (-3.8, 1.25) {\textbf{Microarchitecture}\\[-1pt]\subline{Pipeline $\cdot$ Branch Prediction $\cdot$ Out-of-Order $\cdot$ Superscalar}};
\node[box, fill=c1] (cl1) at (-3.8, 0.05) {\textbf{Hardware}\\[-1pt]\subline{CPU $\cdot$ GPU $\cdot$ FPGA $\cdot$ ASIC $\cdot$ Transistors}};

% Column boxes, right
\node[box, fill=m6] (ml6) at ( 3.8, 6.05) {\textbf{Application Layer}\\[-1pt]\subline{Multi-Agent $\cdot$ Task Decomposition $\cdot$ Cross-Validation $\cdot$ Consensus}};
\node[box, fill=m5] (ml5) at ( 3.8, 4.85) {\textbf{Orchestration Layer}\\[-1pt]\subline{Agent Runtime $\cdot$ MCP $\cdot$ Sandboxing $\cdot$ Planning/Monitoring}};
\node[box, fill=m4] (ml4) at ( 3.8, 3.65) {\textbf{Semantic Interface}\\[-1pt]\subline{Prompting $\cdot$ ICL $\cdot$ Structured Output $\cdot$ Schemas}};
\node[box, fill=m3] (ml3) at ( 3.8, 2.45) {\textbf{Context Management}\\[-1pt]\subline{KV Cache $\cdot$ Context Window $\cdot$ RAG $\cdot$ Attention Patterns}};
\node[box, fill=m2] (ml2) at ( 3.8, 1.25) {\textbf{Inference Serving}\\[-1pt]\subline{vLLM $\cdot$ SGLang $\cdot$ PagedAttention $\cdot$ Speculative Decoding}};
\node[box, fill=m1] (ml1) at ( 3.8, 0.05) {\textbf{Physical Execution}\\[-1pt]\subline{GPU/TPU $\cdot$ ASIC $\cdot$ Neuromorphic $\cdot$ Quantum}};

% Layer numbers aligned at left margin
\foreach \y/\lab in {6.05/L6,4.85/L5,3.65/L4,2.45/L3,1.25/L2,0.05/L1}{
  \node[layernum] at (-6.85,\y) {\lab};
}

% Horizontal guide stubs for layer numbers
\foreach \y in {6.05,4.85,3.65,2.45,1.25,0.05}{
  \draw[bgline,line width=0.45pt] (-6.55,\y) -- (-6.05,\y);
}

% Mapping arrows with relation labels
\draw[rel] (cl6.east) -- node[maplabel] {Scale}     (ml6.west);
\draw[rel] (cl5.east) -- node[maplabel] {Manage}    (ml5.west);
\draw[rel] (cl4.east) -- node[maplabel] {Abstract}  (ml4.west);
\draw[rel] (cl3.east) -- node[maplabel] {Store}     (ml3.west);
\draw[rel] (cl2.east) -- node[maplabel] {Compute}   (ml2.west);
\draw[rel] (cl1.east) -- node[maplabel] {Substrate} (ml1.west);

% Right-side braces / planes
\draw[decorate,decoration={brace,amplitude=5pt},line width=0.9pt,draw=purple!55]
  (6.75,3.05) -- (6.75,0.18)
  node[midway,right=9pt,sidetitle,text=purple!60!black] {Probabilistic\\Execution Plane};

\draw[decorate,decoration={brace,amplitude=5pt},line width=0.9pt,draw=teal!60!black]
  (6.75,6.6) -- (6.75,3.28)
  node[midway,right=9pt,sidetitle,text=teal!60!black] {Deterministic\\Control Plane};

% Substrate indicator for L1
\draw[dashed,draw=softgray!60,line width=0.6pt] (6.75,0.48) -- (6.75,-0.55);
\node[font=\scriptsize,text=softgray] at (7.85,-0.15) {Physical Substrate};

% Small explanatory note
\node[font=\scriptsize,text=softgray,align=center] at (0,-0.72)
{Six-layer analogy from classical systems abstractions to model-native computing components};

\end{tikzpicture}%
}
\caption{Layer comparison between the classical computing stack and the model-native computing stack.}
\label{fig_dual_stack}
\end{figure}

